% =========================================================
% Quellenverzeichnis
%
% Variante A (Standard): alle Quellen in einer Liste
% → einfach \printbibliography verwenden (unten auskommentiert)
%
% Variante B (FHDW-Konvention): Quellen nach Typ getrennt
% → erfordert keyword-Felder in der .bib-Datei:
%     keywords = {mono}   für Monographien
%     keywords = {mag}    für Aufsätze / Zeitschriften
%     keywords = {web}    für Internetquellen
%     keywords = {art}    für Zeitungsartikel
%     keywords = {leg}    für Rechtsprechung
%     keywords = {comp}   für Unternehmensunterlagen
% =========================================================

\clearpage
\phantomsection
\addcontentsline{toc}{section}{Quellenverzeichnis}
\section*{Quellenverzeichnis}

\setlength\bibitemsep{1.5\itemsep}
\setlength{\bibhang}{2em}

% ----- Variante A: eine ungekürzte Liste -----
\printbibliography[heading=none]

% ----- Variante B: nach Quelltyp getrennt (auskommentiert) -----
% \begingroup
% \sloppy
% \defbibheading{mono}{\subsection*{Monographien}}
% \defbibheading{mag}{\subsection*{Aufsätze in Sammelbänden und Zeitschriften}}
% \defbibheading{web}{\subsection*{Internetquellen}}
% \defbibheading{art}{\subsection*{Zeitungsartikel}}
% \defbibheading{leg}{\subsection*{Rechtsprechung}}
% \defbibheading{comp}{\subsection*{Unternehmensunterlagen / Gesprächsnotizen}}
%
% \printbibliography[heading=mono,keyword=mono]
% \printbibliography[heading=mag,keyword=mag]
% \printbibliography[heading=web,keyword=web]
% \printbibliography[heading=art,keyword=art]
% \printbibliography[heading=leg,keyword=leg]
% \printbibliography[heading=comp,keyword=comp]
% \endgroup
